\documentclass[output=paper,colorlinks,citecolor=brown]{langscibook}
\ChapterDOI{10.5281/zenodo.18845604}
\author{Mejrema Koca\orcid{}\affiliation{} and Ulrich Mehlem\orcid{}\affiliation{}}
\title[Perspectives of newly arrived Kurdish and Afghan families in Germany]{Perspectives of newly arrived Kurdish and Afghan families in Germany on their heritage languages and linguistic situation}
\abstract{The linguistic situation of newly arrived families threatened by forced migration in Germany after 2015 is characterised by attempts to learn German as the new official language and an unclear status of other languages, which played a role in their earlier lives, such as the heritage language of the family, the official language of the country of origin and other transit languages \citep{bruecker_rother_schupp2016,panagiotopoulou_uca2023_dynamic,panagiotopoulou_uca_samani2023}. While the status of minority languages without official support in European countries and the implications for integration and educational success are an issue of extensive debate \citep{kristen_et_al2011,kruse_et_al2012,kempert_et_al2016}, the particular cases of Kurds from Syria and Iraq and speakers of Pashto from Afghanistan are rarely studied. On the basis of semiguided interviews, the project examines the development of attitudes of minority school-age children and their parents from Afghanistan, Syria and Iraq on their different languages and their family language policy in a longitudinal perspective. In the theoretical framework of \citet{king_fogle_logan-terry2008} and \citet{darvin_norton2015,darvin_norton2017}, their linguistic resources, preferences, ideologies and investments in language learning are analysed. First project results will be presented on the basis of case studies grouped together in a typology on the basis of \citet{brizic2021}.}

\IfFileExists{../localcommands.tex}{
   \addbibresource{../localbibliography.bib}
   \usepackage{tabularx,multicol}
\usepackage{url}
\urlstyle{same}
\usepackage{textcomp}
 
\usepackage{langsci-optional}
\usepackage{langsci-lgr}
\usepackage{langsci-gb4e}
\usepackage{langsci-branding}

% \usepackage{fontspec}
\babelprovide{georgian} % Georgian language
\babelprovide{armenian} % Armenian language
\babelfont[georgian,armenian]{rm}{DejaVu Sans}
\newfontfamily\georgianfont{DejaVu Sans}
\newfontfamily\armenianfont{DejaVu Sans}

\usepackage{dialogue}
\usepackage{attrib}
\usepackage{attrib} % speaker labels in Koca's chapter
\usepackage{phoenician} % alf, bet in Mengozzi's chapter

   

\makeatletter
\let\thetitle\@title
\let\theauthor\@author
\makeatother

\newcommand{\togglepaper}[1][0]{
   \bibliography{../localbibliography}
   \papernote{\scriptsize\normalfont
     \theauthor.
     \titleTemp.
In Saloumeh Gholami \& Ulrich Mehlem (eds.), \textit{Identity formation, language and migration dynamics: Minorities in the MENA region and the German diaspora}, Berlin: Language Science Press [preliminary page numbering]
   }
   \pagenumbering{roman}
   \setcounter{chapter}{#1}
   \addtocounter{chapter}{-1}
}

\newbool{bookcompile}
\booltrue{bookcompile}
\newcommand{\bookorchapter}[2]{\ifbool{bookcompile}{#1}{#2}}


\newcommand{\georgian}[1]{{\georgianfont #1}}
\newcommand{\armenian}[1]{{\armenianfont #1}}

\newcommand{\textarab}[1]{{\arabicfont \beginR #1} \endR}
\newcommand{\texthebrew}[1]{{\hebrewfont \beginR #1} \endR}
\newcommand{\textsyriac}[1]{{\syriacfont \beginR #1} \endR}


\newfontfamily\phoenicianfont{Tuffy Regular.ttf}

   %% hyphenation points for line breaks
%% Normally, automatic hyphenation in LaTeX is very good
%% If a word is mis-hyphenated, add it to this file
%%
%% add information to TeX file before \begin{document} with:
%% %% hyphenation points for line breaks
%% Normally, automatic hyphenation in LaTeX is very good
%% If a word is mis-hyphenated, add it to this file
%%
%% add information to TeX file before \begin{document} with:
%% %% hyphenation points for line breaks
%% Normally, automatic hyphenation in LaTeX is very good
%% If a word is mis-hyphenated, add it to this file
%%
%% add information to TeX file before \begin{document} with:
%% \include{localhyphenation}
\hyphenation{
    par-a-digm
    Ver-to-vec
    Go-go-lin
    sa-gen
    mo-misch
    ant-wor-te-te
    dix-wa-zim
    spre-chen
    B-Verf-GE
    Swe-dish
    re-mains
    schrei-ben
    di-pey-vin
    Chris-ten-sen
    Wes-sen-dorf
    Ja-va-khi-shvi-li
    ana-lysis
}

\hyphenation{
    par-a-digm
    Ver-to-vec
    Go-go-lin
    sa-gen
    mo-misch
    ant-wor-te-te
    dix-wa-zim
    spre-chen
    B-Verf-GE
    Swe-dish
    re-mains
    schrei-ben
    di-pey-vin
    Chris-ten-sen
    Wes-sen-dorf
    Ja-va-khi-shvi-li
    ana-lysis
}

\hyphenation{
    par-a-digm
    Ver-to-vec
    Go-go-lin
    sa-gen
    mo-misch
    ant-wor-te-te
    dix-wa-zim
    spre-chen
    B-Verf-GE
    Swe-dish
    re-mains
    schrei-ben
    di-pey-vin
    Chris-ten-sen
    Wes-sen-dorf
    Ja-va-khi-shvi-li
    ana-lysis
}

   \boolfalse{bookcompile}
   \togglepaper[6]%%chapternumber
}{}

% \setlength{\footheight}{49.47pt}
\newenvironment{interview}{
  \setlength{\leftskip}{2.5em} % Left margin indent
  \setlength{\rightskip}{2.5em} % Right margin indent
  \setlength{\parindent}{0pt} % Remove paragraph indentation
  \vspace{\topsep} % Add space before, like quote 
}{
  \vspace{\topsep} % Add space after, like quote
}

\newcommand{\speaker}[1]{%
  \noindent\hangindent=2em\makebox[1.5em][l]{#1:}\hspace{0.5em}%
}
\newcommand{\myattrib}[1]{%
    \attrib{#1}
}

\begin{document}
\maketitle

\section{Introduction}\label{ch2.5.ss.1}

In 2015 and 2016, Germany witnessed the influx of more than one million people from the Middle East and the Balkans who fled civil wars in Syria and Iraq, political instability in Afghanistan and persecution of the \isi{minorities} in Turkey, Iran and other countries. According to official data, the most important countries of origin of 441,899 asylum seekers in Germany in 2015 came from Syria (35.9\%), Albania (12.6\%), Afghanistan (7.1\%) and Iraq (6.7\%). In 2016, out of 722,370, 36.9\% were of Syrian, 17.6\% of Afghan and 13.3\% of Iraqi origin \citep{bamf2018}. Among these, there was an overproportional share of children and young adults. Germany's educational system was not well prepared for such a situation. Although the educational success of the' children of migrants who live a long time in Germany remained an issue of debate, the school system had slowly moved from a long tradition of separate classes for learners of \ili{German} as a second language towards the immediate integration of all children into regular classes, supported by additional measures of linguistic improvement. \citep{brueggemann_nicolai2016}. As a consequence of the arrival of many school-age children without knowledge of \ili{German}, the parallel or at least transitional parallel system of special classes witnessed a revival in all the Federal states of Germany. In these classes, the teaching of \ili{German} as a second language, reading, writing, and mathematics was the focus. After a period of one year, maximally two years, the newly arrived children were transferred to regular classes \citep{von_dewitz_massumi2017}. In the aftermath of these reforms, a discussion started among educational researchers, linguists, teachers, and politicians about the best way to integrate these children: some argued in favor of the parallel or partly integrative system, because language teaching can take the specific linguistic and psycho-social conditions of the students directly into account. On the other hand, language learning is not supported by contact with other children of \ili{German} language. The model of complete integration into regular classes has the advantage of more intensive informal language contact, but less structured input \citep{karakayali_zur_nieden2018,roehner2017,roehner2020}.

\largerpage
As an important aspect of this debate, the role of other languages than \ili{German} in the life of the newly arrived children is considered. Many groups of migrants lived in multilingual societies. They do not only speak one language in the family, but distinguish between the language of a local ethnic community and the official language of the country of origin. A survey of newly arrived migrants in 2016 \citep{bruecker_rother_schupp2016} showed that 76\% of Afghanistan refugees spoke \ili{Farsi} and 20\% \ili{Pashto}, 78\% of Syrians spoke \ili{Arabic} and 20\% Kurmanji\il{Kurdish!Kurmanji}, 50\% of Iraqis spoke \ili{Arabic}, 37\% Kurmanji\il{Kurdish!Kurmanji} and 9\% Sorani. Languages of third countries where the children lived during their flight may even play a role.

The project on ``Linguistic integration of newly arrived primary school children in Germany'' as part of the LOEWE research cluster on ``Minority Studies: Language and \isi{Identity}'' started with the following research questions: 

\begin{enumerate}
    \item What are the language attitudes\is{Language!attitudes} of newly arrived migrant families towards their different languages and \ili{German}? 
    \item What kind of language policy adopt the families during the first year of their stay? 
\end{enumerate}

To answer these questions, we conducted semistructured interviews with children and their parents about their language attitudes\is{Language!attitudes} and orientations and their situation in Germany immediately after their enrollment in the school and 18 months later. Additionally, the children's knowledge of \ili{German} and the heritage language was tested three times with a picture description test. The children spoke Kurmanji\il{Kurdish!Kurmanji} \ili{Kurdish} (3), \ili{Pashto} (3) and \ili{Dari}/\ili{Farsi} (9) as their first languages. Their families came from Syria, Iraq, Iran, and Afghanistan. In the following, we will present some results of interview analysis of \ili{Kurdish} and \ili{Pashto}-speaking children and their parents.

The paper consists of five sections: In Section \ref{ch2.5.ss.2}, the theoretical basis of the research will be presented. The role of first language(s) in migration contexts, the specific conditions of \isi{minorities} and newly arrived families and the model of ideology, identity and investment will be discussed. Section \ref{ch2.5.ss.3} deals with the methodology of the interviews, especially with children. Section \ref{ch2.5.ss.4} presents an analysis of 10 interviews with a focus on language attitudes\is{Language!attitudes} of parents and children at two different time points. The data will be discussed in a comparative perspective on the basis of Norton's model. A first attempt will also be made to establish a typology that is credited to the work of \citet{brizic2021}. Section \ref{ch2.5.ss.5} will discuss the results and draw some conclusions.

\section{Theoretical background: The role of languages brought to Germany}\label{ch2.5.ss.2}

\subsection{Heritage, official, literacy and transit languages}\label{ch2.5.ss.2.1}

In the literature on migration and multilingualism, several terms are used to define a language acquired before the onset of second language acquisition\is{Language!acquisition} in the immigration country:
\newpage
\begin{itemize}
    \item first language (focus on the time of acquisition\is{Language!acquisition}), 
    \item mother tongue (focus on the person of basic relationship), 
    \item family language (focus on the role of the family for the transmission of a language) 
    \item language of origin (focus on a starting point of an individual's or communities life)
    \item heritage language (focus on conveying a language from ancestors to descendents), 
\end{itemize}
    
We will use \textit{heritage language} in the following. The term refers to the first language acquired by an individual who grows up in a family in which the language of the surrounding majority society is not (exclusively) used \citep[131]{polinsky2018}. Such a constellation also exists in the context of migration -- regardless of how far back in time the person arrived. 

In the case of linguistic/ethnic \isi{minorities} a tension between the heritage language and the official language of the country of origin can be observed. In addition, there might be other languages used in other communities living in close contact to the own family. At school, only the official language might be taught as a literacy language. During the flight and prior to the arrival in Germany, the children might have been in contact with languages of other countries: transit languages \citep{le_pichon-vorstmann2020}. In order to understand the complexity of such forms of multilingualism, the following paper documents how -- in the sense of sociolinguistics and linguistic anthropology -- parents and children perceive and evaluate their linguistic practice. 

\subsection{Importance of heritage languages in theories of integration/acculturation}\label{ch2.5.ss.2.2}

In the processes of migration, acculturation and integration, heritage languages -- at the side of the official language of the immigration country -- play an important role. As international studies on educational success such as PIRLS 2016 \citep{mullis_et_al2017} show, these languages are still spoken in many families at least sometimes (14\%), almost always (17\%) or exclusively (5\%). In Germany, the corresponding figures are 16\%, 15\% and 1\%. On the basis of PISA data for 2012, \citet{gebhardt_rauch_mang_saelzer_stanat2013}, confirm a usage of L1 for students in grade 9 of 64\% in first generation (migrated by themselves) and 34\% in second generation (born as children of migrants). 

Studies on language attitudes\is{Language!attitudes} \citep{baker1992} present data on a certain adherence to heritage languages, even for the younger generation \citep{boos-nuenning_karakasoglu2015,broeder_extra1999}, in migration contexts. These attitudes\is{Language!attitudes} are part of what is called family language policy in more recent studies \citep{king_fogle_logan-terry2008}.

\isi{Heritage languages} may develop differently than in the context of origin, but still reach a high level of proficiency\is{Language!proficiency} among immigrant groups. Although a strict causality between a high competence in the first and the second language, as it was claimed by the interdependence hypothesis \citep{cummins1979}, does not exist, a good knowledge of the first language does not exclude that of the second. For children of Russian, \ili{Turkish} and Vietnamese origin in Germany, \citet{knigge_klinger_schnoor_gogolin2015} found a small, albeit significant influence of L1 proficiency\is{Language!proficiency} on L2 proficiency\is{Language!proficiency}: 6\% of variance of the latter could be explained by the first language competence. 

According to the theory of segmented assimilation \citep{portes_fernandez-kelly_haller2009,kristen_olczyk2013}, the maintenance\is{Language!maintenance} of the first language may even help in developing a profile of double integration, with respect to the greater society of the immigration country and the community of origin in the categories of Berry's acculturation theory \citep{berry1997,berry_phinney_sam_vedder2006}. In terms of national background, the data of \citet{knigge_klinger_schnoor_gogolin2015} showed an advantage in L2 proficiency\is{Language!proficiency} in the Vietnamese speakers, with speakers of Russian at the second and speakers of \ili{Turkish} at the third place, even after control for socio-economic background.

\subsection{Research on heritage languages of minorities in migration contexts}\label{ch2.5.ss.2.3}

The research presented so far refers to migrant groups and heritage languages that are at the same time the official languages of the country of origin. This is not the case for the ethnic and linguistic \isi{minorities} discussed here: The \ili{Kurdish} variety of Kurmanji\il{Kurdish!Kurmanji} is spoken in Turkey, Northern Syria and Iraq, but is not an official language in Turkey and Syria and only accepted in the autonomous region of Kurdistan in Northern Iraq. \ili{Pashto} in Afghanistan, although constituting one of the two official languages, is of minor importance in provinces where the majority speaks other languages and uses \ili{Dari} as \textit{lingua franca}. In Germany, these languages are part of official statistics, but still play only a minor role in teaching \citep{brehmer_mehlhorn2018}.

\citet{brizic2006,brizic_hufnagl2016,brizic_bulut_simsek2021,brizic2021} demonstrate for \ili{Kurdish} migrants from Turkey in Austria how processes of language maintenance\is{Language!maintenance} and/or shift of the generation of the parents in Turkey have an impact on the educational success of the children in the migration country. With respect to narrative competence in \ili{German} however, multilingual and bilingual children do not differ significantly: Maintenance of the first language \ili{Kurdish} or language shift to the official language \ili{Turkish} does not lead to differences in narrative competence in \ili{German} \citep[215--228]{brizic2021}.

The Moroccan immigrants of \ili{Berber} origin in Germany \citep{mehlem1998} and the Netherlands \citep{leseman_scheele_mayo_messer2009} also constitute an interesting reference point of research. The restricted linguistic input in \ili{Berber}, measured in a parent questionnaire on language and media use at home, leads to a slower growth in vocabulary and text comprehension, compared to children of \ili{Turkish} origin. On the contrary, the development of \ili{Dutch} vocabulary and morphology is more successful in the \ili{Berber}-speaking group, perhaps due to a greater input in \ili{Dutch}. \citep[303]{leseman_scheele_mayo_messer2009}. But in both groups, after controlling for \ili{Dutch} vocabulary knowledge, short-term memory and Raven IQ at age 3, the variance of \ili{Dutch} proficiency\is{Language!proficiency} at the age of six years is also explained to some extent by the oral proficiency\is{Language!proficiency} in the first language at age 3 (albeit with small effect sizes, \citealt[308]{leseman_scheele_mayo_messer2009}.

As the case studies on speakers of \ili{Kurdish} and Tarifit\il{Berber!Tarifit} \ili{Berber} show, \isi{minority} languages already in the country of origin have specific effects on linguistic and cultural development in migration contexts. However, the results still remain contradictory. 

Most of the studies mentioned so far, even if they look at \isi{minorities} and their languages, deal with migrants families who have been living already for a long time in the country of immigration. The results in linguistic development, language attitudes\is{Language!attitudes} and use cannot be transferred to the situation of newly arrived migrants who constitute a growing part of the migrant population. 

\subsection{Language ideology\is{Language!Ideology}, identity and investment}\label{ch2.5.ss.2.4}

Language attitudes\is{Language!attitudes} of parents and children are deeply influenced by language ideologies as sets of ideas attributing prestige and values to languages and linguistic varieties. They are ``… dominant ways of thinking that organize and stabilize societies while simultaneously determining modes of inclusion and exclusion'' \citep[72]{darvin_norton2015}.

These ideas circulate in a given society and are often considered as the natural order of things. In such a way, there is strong pressure in a country like Germany to learn the national language quickly. It is postulated that the \ili{German} language is the ``key to integration'' that only leads to acceptance and success \citep{gogolin1994}. On the other hand, a heritage language can also play a role in an imagined community as a value in itself to be defended against outside pressure. Finally, a national language such as \ili{Arabic} in Syria or \ili{Turkish} in Turkey may be identified with legitimacy and ethnic purity. 

Language ideologies establish a link between languages and group identities which may play a major or minor role in the development of a child's personal identity. In the life cycle, an individual has to solve different tasks such as the choice of a profession, the foundation of a new family or the adherence to a greater or smaller collective. The choice of a language might play an important role in the shaping of identity.

\citet{darvin_norton2015} introduce the term ``capital'' according to \citegen{bourdieu1986} theory to identify economic (income, financial resources, professional status), social (networks, relationships) and cultural resources (educational qualifications, cultural assets, cultural skills, including language skills) an individual can mobilize.

The last element of the model, investment, \citet{darvin_norton2015} describes all kinds of efforts undertaken by an individual in order to achieve a specific goal. The notion of investment ``… accentuates the role of human agency and identity in engaging with the task at hand, in accumulating economic and symbolic capital, in having stakes in the endeavor and in persevering in that endeavor'' \citep{kramsch2013}.

\citet{darvin_norton2015} recognizes that, under the conditions of globalization, the spaces in which language acquisition\is{Language!acquisition} and socialization take place are increasingly ``deterritorialized and unbounded''. This also holds for the case studies that will be discussed in the following. 

\section{Qualitative interviews with children and parents: An intergenerational perspective}\label{ch2.5.ss.3}

In order to answer the research questions on language attitudes\is{Language!attitudes} (1) and the family language policy (2) of the newly arrived children and their parents, semistructured interviews were conducted on the basis of questions dealing with language practices in the family and school experience in the home country, contact with other languages in transit countries and the situation at the beginning in Germany, leading to expectations for the future.

Before delving further into the research design, the question, how to design and conduct interviews in a way appropriate for children has to be clarified. Several studies already address exactly these questions, namely how qualitative interviews especially with children can or should be designed. It has been found that some common methods of qualitative interview design, which were implicitly designed for adult research subjects, do not seem to be suitable for children. Against this background, this study chose a child-specific strategy for conducting interviews, which is also reflected in the design of the guideline. These strategies are mainly based on ideas from \citet{delfos2008}.

There is a tension between a structured interview strategy and a flexibility that allows children to bring in personal views and interpretations. It should be found out, which issues the interviewed persons themselves would like to emphasize \citep{delfos2008}. Therefore, a mixed form of guided and partially narrative interview was chosen \citep{helfferich2009,delfos2008}. Data analysis is carried out with the help of qualitative content analysis according to \citet{mayring2010}.

The interviews in \ili{German} were conducted by M. Koca. She was also present during the interviews in \ili{Kurdish} and \ili{Pashto} which were administered by field assistants speaking the heritage language of the families on the basis of translated guidelines. The children were interviewed first and in the absence of the parents but with their prior written consent. Parent interviews followed later. If the family did not accept an interpreter from outside, older children occasionally served as translators during the conversations with parents. Such an approach enabled to consider both perspectives independently from one another.

Ethical issues of interviewing particularly vulnerable young children because of their experience of forced migration also were at stake. The project was approved by the Ethics Board of the Department of Education of Goethe University Frankfurt. The careful preparation of all participants and the transparent communication of data protection regulations also contributed to the quality of the interviews. The explanation and signing of a data protection declaration were not only of legal necessity, but also offered the interviewees emotional security and encouraged an atmosphere of trust during the interviews.

The researchers themselves do not claim to leave their own experiences in the field completely outside: One of the researchers had experienced forced migration herself. In her situation, the uncertainty about whether she would be permitted to remain in Germany was a greater burden than the pressure to learn the language, even though she was able to study. The other researcher, engaged in teacher education and the strengthening of multicultural \isi{attitudes} over a long time period had to critically reflect his approach of defending linguistic diversity as a value in itself.

\section{Interview analysis: Five case studies}\label{ch2.5.ss.4}

\subsection{Data collection and overview of the sample}\label{ch2.5.ss.4.1}

According to the aim of the project to focus on newly arrived migrants with primary school children, originating from ethnic and linguistic \isi{minorities} of the Middle East, it was planned to contact the families via the \ili{German} schools directly after the enrollment of the children. Due to the Covid pandemic and the closure of schools for long periods in 2020 and 2021, the first families were contacted in 2020 via an organisation running a hostel for refugees. They belong to \ili{Kurdish}-speaking \isi{minorities} from Syria and Iraq and had come to Germany as early as 2015 and 2016 (cases 1 and 3). 

Only in 2021, a direct cooperation with two primary schools was possible. These schools were running \textit{Willkommensklassen} `welcome classes' according to the transitional parallel model of teaching in the \ili{German} state of Hesse. 13 children and their families were contacted. These children spoke \ili{Pashto} (3) and \ili{Dari}/\ili{Farsi} (9) and Kurmanji\il{Kurdish!Kurmanji} \ili{Kurdish} (1) as their first languages. Their families came from Syria, Iran and Afghanistan. Out of these, only five cases are presented in this article. At the time of the first interview, they were living in Germany between 7 months and 1.5 years. The following table gives some basic information about the five cases and their languages.

\begin{table}%[htp]
    %\centering
    \fittable{\begin{tabular}{lllll}
    \lsptoprule
& Arrival in & & Heritage and &\\
& Germany & \is{Origin!country}Origin Country & Transit Language& \isi{Literacy} \\
\midrule
Z (15 years) & October 2015 & Syria (\ili{Kurdish} Region, & \ili{Kurdish} (Kurmanji\il{Kurdish!Kurmanji}), & \ili{Turkish}, \\ 
Case 1 & & Northern Syria) & \ili{Turkish}, \ili{Arabic} & \ili{German} \\
\midrule
N (10 years) & March 2021 & Syria (\ili{Kurdish} Region, & \ili{Kurdish} (Kurmanji\il{Kurdish!Kurmanji}), & \ili{German} \\
Case 2 & & Northern Syria) & \ili{Arabic} & \\
\midrule
A (13 years) & February 2017 & Iraq (Sinjar province?) & \ili{Kurdish} (Kurmanji\il{Kurdish!Kurmanji}), & \ili{Kurdish}, \\
Case 3 & & & Variety Bahdini\il{Kurdish!Bahdini} & \ili{German} \\
\midrule
S (9 years) & February 2020 & Afghanistan & \ili{Pashto}, \ili{Urdu} & \ili{German} \\
Case 4 & & & \\
\midrule
E (8 years) & February 2020 & Afghanistan & \ili{Pashto}, \ili{Farsi} & \ili{German} \\
Case 5 & & & \\
\lspbottomrule
    \end{tabular}}
    \caption{Overview of the case studies}
    \label{ch2.5.tab.1}
\end{table}

Two of the three \ili{Kurdish} (Kurmanji\il{Kurdish!Kurmanji}) speaking families originate from Syria, and one from Iraq. The two Afghan families lived in the provinces of Baghlan and Jalalabad in Afghanistan and spoke \ili{Pashto} as their first language.  As transit languages, \ili{Turkish} (case 1: stay in Turkey) and \ili{Urdu} (case 4: stay in Pakistan) play a role. 

The interview analysis is based on two categories: ``Reported use and competence of languages brought to Germany'' (\sectref{ch2.5.ss.4.2}) gives information on language use and language competence in the family and in the languages learnt at school (literacy). The second category ``Cultural identity and attitudes towards family languages towards family languages'' (\sectref{ch2.5.ss.4.3}) discusses questions of family language policy: preferences, maintenance\is{Language!maintenance} and shift.

The following excerpts from the interview transcripts contain passages in different languages according to the choice of the interviewees.\footnote{I = interviewer, MA = A's mother, FN = N's father, MN = N's mother, MS = S's mother, FE = E's father, MZ = Z's mother.} Sometimes, statements of children and parents are presented together. In two parent interviews, the children also translated from \ili{Kurdish} to \ili{German} and vice versa. These translations were part of the interview itself. The author's translations in \ili{English} are always given in single quotes '…'. If (translates) is written after the object language, the speakers were immediately translating their utterances for the other party. All utterances were transcribed verbatim and are not corrected for the standard language. The excerpts presented in this article include data of the first and second interviews (18 months later).   

\subsection{Reported use and competence of languages brought to Germany}\label{ch2.5.ss.4.2}

\subsubsection{Z: Kurdish in Syria, transit via Turkey}\label{ch2.5.ss.4.2.1}

Z and her mother and sisters fled from Syria. They speak \ili{Kurdish} at home. Before travelling to Germany, they spent three years in Turkey, where they all learnt \ili{Turkish}. In the following, the child describes her use of different languages at home.

\begin{interview}\label{ch2.5.ex.1}
\speaker{I}\textit{In welchen Sprachen sprichst du zu Hause?} \\
`Which languages do you speak at home?'  

\speaker{Z}\textit{Ähm, Kurdisch und wenn meine Geschwister mich, also manche Wörter nicht verstehen, dann erklär' ich die auf Türkisch.} \\
`\ili{Kurdish}, and if my sister and brother don't understand me, well, some words, then I explain them in \ili{Turkish}.' 

\speaker{I}\textit{Mhm, und warum auf Türkisch?} \\
`Mhm, and why in \ili{Turkish}?' 

\speaker{Z}\textit{Weil, die können nicht so viel Kurdisch, also sie sprechen schon Kurdisch, aber nicht so.} \\
`Because they don't speak so much \ili{Kurdish}, so they already speak \ili{Kurdish}, but not enough.' 
\myattrib{Int 1\_Z I, 21--26}
\end{interview}

\ili{Turkish} is used here as a language of mutual understanding between Z and her younger siblings. Due to the long stay in Turkey, they also used \ili{Turkish} at home. Since their arrival in Germany, the knowledge of \ili{Kurdish} decreased, at least in the case of the younger siblings. With respect to literacy, the interviewer asks Z about her school experience in Syria. 

\begin{interview}\label{ch2.5.ex.2}
\speaker{I}\textit{So liebe Z, ich möchte dich kurz fragen, in welcher Sprache kannst du schreiben?} \\
`Dear Z, I would like to ask you, in which language can you write?' 

\speaker{Z}\textit{Ähm Türkisch und Deutsch} \\
`Um, \ili{Turkish} and \ili{German}' 

\speaker{I}\textit{Und Arabisch?} \\
`And \ili{Arabic}?' 

\speaker{Z}\textit{Nein kann ich nicht} \\
`No I can't' 

\speaker{I}\textit{Und du warst in der Schule in in Syrien?} \\
`And you went to school in Syria?' 

\speaker{Z}\textit{Ja hm nur äh bis zweite Klasse isch kann nur manchmal so etwas lesen und kann nur mein Name schreiben sonst nicht.} \\
`Yes hm only uh until second grade I can only sometimes read something like and can only write my name otherwise not.' 

\speaker{I}\textit{Und habt ihr dort in der Schule Arabisch gelernt oder welche Sprache gab's da in der Schule?} \\
`And did you learn \ili{Arabic} there in school or what language was there in school?' 

\speaker{Z}\textit{Arabisch und als äh also und Englisch so wie hier} \\
`\ili{Arabic} and as uh so and \ili{English} so like here' 

\speaker{I}\textit{Hm, aber Arabisch hast du dort gelernt zu schreiben?} \\
`Hm, but \ili{Arabic} you learned there to write?' 

\speaker{Z}\textit{Ja} \\
`Yes' 

\speaker{I}\textit{Und jetzt kannst du nicht mehr oder wie ist es?} \\
`And now you can't write it or how is it?' 

\speaker{Z}\textit{Äh ich habe vergessen, weil ich war ja in drei Jahre in Türkei… und hier sieben Jahre} \\
`I forgot, because I was in Turkey for three years… And here for 7 years' 
\myattrib{Int 1\_Z II, 1--30}
\end{interview}

Z claims \ili{Turkish} and \ili{German} as her languages of literacy. Only after directly asked, she speaks about \ili{Arabic} which was the first official school language in Syria. However, the amount of time spent in school was too short, and later, \ili{Turkish} and \ili{German} replaced it completely as written language. 

\begin{interview}\label{ch2.5.ex.3}
\speaker{I}\textit{Ah ok. Und in der Türkei, wie hast du dort hm Türkisch gelernt zu schreiben?} \\
`And in Turkey, how did you learn to write in \ili{Turkish}?' 

\speaker{Z}\textit{Hm, mit Leute die Türken waren, \textnormal{(I: \textit{hm})} die waren befreundet und ja von den haben wir was gelernt und die dürften in so eine Moschee gehen immer, so in Sommerferien 6 Monate lang und dort habe ich auch gelernt Türkisch zu sprechen. Aber schreiben und lesen habe ich hier gelernt nach Deutschen Buchstaben und so.} \\
`With people who were \ili{Turkish}. They were friends and we learnt from them and they were allowed to go to the mosque, in summer holidays for 6 months and there I learnt to speak \ili{Turkish}. But reading and writing \ili{Turkish} with \ili{German} letters I learnt here.' 
\myattrib{Int 1\_Z II, 32--38}
\end{interview}

It is not clear which form of language lessons the children visited in Turkey, but it was \ili{Turkish}, not \ili{Arabic} which was taught there. Interestingly, only the spoken language was practiced there.\footnote{According to \citet[56]{chiari2018} such informal teaching for Syrian refugees falls in the first phase of educational policy (2011--2014), in which the \ili{Turkish} state did not systematically enroll Syrian refugee children in the educational system.} In Germany, \ili{Turkish} still serves as a bridge to acquire new \ili{German} words, when the girl uses Google Translate to find corresponding expressions to \ili{Turkish}. \ili{Kurdish} is not mentioned at all in this part of the interview.

In the second interview, Z was asked which language she felt most comfortable with (\textit{Wohlfühlsprache}). Her answer is as follows:

\begin{interview}\label{ch2.5.ex.4}
\speaker{I}\textit{In welcher Sprache fühlst du dich wohl, oder welche Sprache sprichst du am liebsten?} \\
`In which language do you feel most comfortable? Which language do you like to speak most?' 

\speaker{Z}\textit{äh Türkisch} \\
`\ili{Turkish}' 

\speaker{I}\textit{Türkisch?} \\
`\ili{Turkish}?' 

\speaker{MA}\textit{und Kurdisch auch beide} `and \ili{Kurdish}, both'

\speaker{I}\textit{Ah ok. Hat sich Türkisch bei dir noch so festgehalten wie damals.} \\
`You kept \ili{Turkish} from that time.' 

\speaker{Z}\textit{Ja} (smiling) \\
`Yes' 

\speaker{I}\textit{Ah schön. Wo hörst oder sprichst du noch mit jemand Türkisch aus der Familie?} \\
`That's nice. Where do you hear \ili{Turkish}, does somebody speak with you \ili{Turkish} in your family?' 

\speaker{Z}\textit{Ja mit Freunde aus der Schule, mit Geschwister, Cousine usw. auch} \\
`Yes, with friends at school, with siblings, cousins, too' 
\myattrib{Int 2\_Z, 36--49}
\end{interview}

For the girl, \ili{Kurdish} is not exclusively the language to express emotions. The use of \ili{Turkish} is not restricted to the greater family, but also includes other children from the new school. Later, she adds that her contact with these children is restricted to school.
 
\subsubsection{N: Kurdish in Syria, transit via Greece}\label{ch2.5.ss.4.2.2}

In case 2, the child claims that Kurmanji\il{Kurdish!Kurmanji}-\ili{Kurdish} is the family language even in Germany. In a sentence, mixed with some \ili{German}, \ili{Kurdish} is also linked to the origin:

\begin{interview}\label{ch2.5.ex.5}
\speaker{I}\textit{Tu xwisk û brê te bi kîjan ziman qise dikî?} \\
`In which language do you speak at home?' 

\speaker{N}\textit{Kurdî} \\
`\ili{Kurdish}…' 

\speaker{I}\textit{Kurmancî-Kurdî?} \\
`Kurmanji Kurdish?' 

\speaker{N}\textit{Kurmanji ist \textnormal{(\ili{German}: is)} zimanê me ye… Em bi eslê xwe kurden \textnormal{(\ili{German}: Kurds}.} \\
`\ili{Kurdish} is our language. We have \ili{Kurdish} origin.' 
\myattrib{Int 1\_N, 12--16, 25--29}
\end{interview}

N learned some \ili{Arabic} only during her long stay in Greece (6 years):

\begin{interview}\label{ch2.5.ex.6}
\speaker{I}\textit{Zimanê din jî dizanî?} \\
`Do you know another language?' 

\speaker{N}\textit{Ez erebî zanim, zimanê xwe û etwas deutsch.} \\
`I know \ili{Arabic}, my native language and a little bit of \ili{German}.' 

\speaker{I}\textit{Tu erebî kîderê hîn buyî?} \\
`Where did you learn \ili{Arabic}? 

\speaker{N}\textit{Li Yunan.} \\
`In Greece.' 
\myattrib{Int 1\_N, 17--24}
\end{interview}

Both parents confirm that \ili{Kurdish} is their mother tongue (\textit{zimanê dayîka}). As for \ili{Arabic}, the father gives the following explanation:

\begin{interview}\label{ch2.5.ex.7}
\speaker{I}\textit{Pistî Kurmancî we kîjan ziman qise dikir?} \\
`Which language do you speak in addition to \ili{Kurdish}?' 

\speaker{FN}\textit{Kurmancî û Erebî. Ne zidetir. Ema zimanê me Kurmancî ye. Em texrîben 5, 6 salî bûn me Erebî destpêkir.} \\
`\ili{Kurdish} and \ili{Arabic}. Not much. But our language is \ili{Kurdish}. We were about 5, 6 years old when we started speaking \ili{Arabic}.' 

\speaker{MN}\textit{Ne zêde.} \\
`Not much.' 
\myattrib{Int 1\_N\_parents, 144--154}
\end{interview}

In the family, \ili{Kurdish} is predominantly used, but also \ili{German} already plays a certain role, at least for the older children. On the one hand, the parents want the children to speak \ili{Kurdish}, on the other hand, they can learn some \ili{German} from them: 

\begin{interview}\label{ch2.5.ex.8}
\speaker{I}\textit{Hûn mal de kîjan ziman dipeyvin?} \\
`Which language do you speak at home?' 

\speaker{FN}\textit{Zarokên min bi Kurmancî dipeyvin. Em jî wan ra Kurmancî dipeyvin. Em wan ra dibêjin Kurmancî bisteqilin, bila em jî wan bigirin. Yên mezin disteqilin. Min ra jî Almanî disteqilin. Ez wan ra dibêjim wele ez fahm nakim. Dibêjin em dixwazin tû fehm bikî. Yên ke hebin bisteqilin, em hîn dibin.} \\
`My children speak \ili{Kurdish}. We speak \ili{Kurdish} with them. We tell them that they should speak \ili{Kurdish}. We learn from them. The older children speak. They speak \ili{German} to me. I say to them that I don't understand. They say that I have to understand/learn it. If people spoke, we would learn from them.' 
\myattrib{Int 1\_N\_parents, 235--240}
\end{interview}

The father also reports on the knowledge of other languages of his children:

\begin{interview}\label{ch2.5.ex.9}
\speaker{I}\textit{Zarokên we wekî din kîjan ziman dipeyvin?} \\
`Which other languages do your children speak?' 

\speaker{FN}\textit{Kurmancî dipeyvin. Û Almanî jî dipeyvin. Yê min ê mezin. 4 hebin terin medresê. Ew xwe jî diştexilin. Kecika min a vir jî disteqile. Kurmancî û Almanî. Zimanê Erebî jî dizanin. Ma ne geleke. Fehm dikin.} \\
`They speak \ili{Kurmanji} Kurdish. They also speak \ili{German}. My daughter, who is here, speaks \ili{Kurdish} and \ili{German}. She can also speak \ili{Arabic}. Isn't that enough?' 
\myattrib{Int 1\_N\_parents, 163--173}
\end{interview}

While Z's family was in Turkey, N's family also stayed a longer time in Greece. But their legal status was not secure. Even worse, the refugee camp was far removed from inhabitated areas:

\begin{interview}\label{ch2.5.ex.10}
\speaker{FN}\textit{Li Yunanê jî em avitin colê, kampa colê de bû. Tistek dora me tune bû. Yani car -- pênc sala zarok necûn medresakî tistekî. Li colê bû.} \\
`In Greece, they threw us into a desert. The camp was in the desert. We had nothing around us. The children couldn't go to school for four to five years. It was in the desert.' 
\myattrib{Int 1\_N\_parents, 348--350, 360--62}
\end{interview}

They could not get in contact with other people outside the camp. Although the older children were already in school age, they could not go to school. It is not surprising that \ili{Greek} language did not play any role in their life. By contrast, the situation in Germany is described quite positively: 

\begin{interview}\label{ch2.5.ex.11}
\speaker{FN}\textit{Li vî derê texrîben 6 heyvê wan terin medresê, ji xwe re dixûnin, dinivîsînin.} \\
`They started school here about 6 months ago. They can read, write.' 
\myattrib{Int 1\_N\_parents, 350--351, 363--64}
\end{interview}

In the second interview, the girl still prefers \ili{Kurdish} as a language she feels comfortable with (Int 2\_N, 29--30). In the new school, however, \ili{German} is predominantly used, because her friends come from several countries. She does not speak \ili{Arabic} with other children: 

\begin{interview}\label{ch2.5.ex.12}
\speaker{N}\textit{Ja, ich kann auch Arabisch, aber ich rede mit denen nicht Arabisch, weil wir sind nicht Araber.} \\
`Yes, I also speak \ili{Arabic}, but I don't speak to them \ili{Arabic}, because we aren't Arabs.' 
\myattrib{Int 2\_N, 55--76}
\end{interview}

The father emphasizes the importance of going to school in Germany and learning. He and his wife want to make some progress, too:

\begin{interview}\label{ch2.5.ex.13}
\speaker{FN}\textit{Em gelekî kêfxweş in bi vê çikê, bo zarokê me jî ku diçin mektebê û ji xwe re dielimin û  \textnormal{(I: \textit{Hm} (agrees))} äh em bi xwe, dixwazim em pêşkevin, ne wek xwe bin.} \\
`We are very happy about it, that our children can go to school and learn something. We want to develop ourselves further, we do not want to stay what we were before.' 
\myattrib{Int 2\_N\_father, 536--537}
\end{interview}

\subsubsection{A: Kurdish in Northern Iraq}\label{ch2.5.ss.4.2.3}

The child originates from the \ili{Kurdish} part of Iraq, from where she and her mother and sisters fled five years ago. Her father had arrived in Germany shortly before. She still speaks \ili{Kurdish}\footnote{According to the \ili{Kurdish} field assistant, it is the Bahdini\il{Kurdish!Bahdini} (also known as Bedini, Behdini, Wehdini) variety of Kurmanji\il{Kurdish!Kurmanji} \citep[S. 287]{haig2019}.} in the family. Her mother confirms her own perspective on a good competence in \ili{Kurdish}, although meanwhile proficiency\is{Language!proficiency} in \ili{German} is the best.

\begin{interview}\label{ch2.5.ex.14}
\speaker{I}\textit{In welchen Sprachen sprichst du zu Hause?} \\
`Which languages do you speak at home?' 

\speaker{A}\textit{Kurdisch} \\
`\ili{Kurdish}' 

\speaker{I}\textit{Kannst du noch eine andere Sprache oder mehrere Sprachen sprechen?} \\
`Can you speak another language or several languages?' 

\speaker{A}\textit{ahm, ich spreche ich kann Deutsch, und Englisch lernen wir gerade in der Schule} \\
`I speak, I can speak \ili{German}, and we are learning \ili{English} at school right now' 
\myattrib{Int 1\_A, 7--11}

\speaker{I}\textit{Hm. Welche Sprache sprichst du mit deinen Geschwistern?} \\
`Which language do you speak with your sisters and brothers?' 

\speaker{A}\textit{ahm: Kurdisch eigentlich meistens} \\
`\ili{Kurdish} for the most time in fact' 
\myattrib{Int 1\_A, 12--13}
\end{interview}

With respect to literacy learning, A refers to her school experiences in Iraq: 

\begin{interview}\label{ch2.5.ex.15}
\speaker{I}\textit{Kannst du schon schreiben oder lesen in einer anderen Sprache?} \\
`Can you already write or read in another language?' 

\speaker{A}\textit{Früher konnte ich Kurdich, aber jetzt habe ich es verlernt, weil ich muss so viel Deutsch lernen. Gerade in Unterricht Englisch} \\
`I was able to read \ili{Kurdish}, but now I forgot it, because I have to learn so much \ili{German}. And now, lessons in \ili{English}' 
\myattrib{Int 1\_A 30--31}

\speaker{I}\textit{Also Deutsch, Englisch kannst du jetzt schreiben} \\
`So, \ili{German}, \ili{English} you are able to write' 

\speaker{I}\textit{Ah. Bist du dort zur Schule gegangen?} \\
`Ah. So did you go to school there?' 

\speaker{A}\textit{Ja, nur bis erste Klasse dann eh muss ich nach Deutschland kommen} \\
`Yes, only until the first grade, then I must come to Germany' 
\myattrib{Int 1\_A, 48--49}
\end{interview}

A was in contact with written \ili{Kurdish} only during a short time, when she went to school in Northern Iraq. The expression \textit{verlernen} `forget' may indicate that \ili{German} as a second language and \ili{English} as the first foreign language replaced \ili{Kurdish} as the first language of literacy.

A's mother also does not refer to \ili{Kurdish} when asked about picture book reading at home. In the following excerpt, A translates the interviewer's question into \ili{Kurdish}, while A's mother answers in \ili{Kurdish}:

\begin{interview}\label{ch2.5.ex.16}
\speaker{I}\textit{Und lesen Sie Ihrem Kind vor oder schauen Sie mit ihm die Bilderbücher an?} \\
`And are you reading to your child or looking at the picture books with them?' 

\speaker{A}\textit{Tu ji zarokên xwe re pirtûkan dixwînî?} (translates)\\
`Do you read books for your children?' 

\speaker{MA}\textit{Nein, manchmal lesen, schreiben} \\
`No, sometimes read write' 

\speaker{I}\textit{In welcher Sprache?} \\
`In which language?' 

\speaker{MA}\textit{Deutsch} \\
`\ili{German}' 

\speaker{I}\textit{Auf Deutsch?} \\
`In \ili{German}?' 

\speaker{MA}\textit{Aber ich ein bisschen schreiben und lesen Deutsch} \\
`But I, a little bit, write and read \ili{German}' 

\speaker{I}\textit{Und auf Kurdisch lessen sie nicht vor?} \\
`And in \ili{Kurdish} you do not read to them?) 

\speaker{MA}\textit{Meine Kind verstehen nicht} \\
`My child does not understand' 
\myattrib{Int 1\_A\_mother, 192--200}
\end{interview}

Picture books are read only in \ili{German}, because the children (A's younger siblings) do not understand enough \ili{Kurdish}: an indication of the withdrawal of \ili{Kurdish} as a family language even in this case. Unfortunately, A and her mother did not participate in the second interview.

\subsubsection{S: Pashto}\label{ch2.5.ss.4.2.4}

When interviewed for the first time, the girl S was nine years old. She was born in Afghanistan and speaks \ili{Pashto} in her family, but she did not learn to write it. 

\begin{interview}\label{ch2.5.ex.17}
\speaker{I}\textit{Ḫa, wadrega. Pe kor ke pe kuma žeba ḫabere kawey?} \\
`Which language do you speak at home?' 

\speaker{S}\textit{Paḫto bānde.} (translates)\\
`In \ili{Pashto}.' 

\speaker{I}\textit{Ḫa, pa kum žebe ke likel ao lwastel kawelaiṣey?} \\
`Which language can you read and write?' 
\myattrib{Int 1\_S, 61ff}

\speaker{I}\textit{Yani, da paḫto ḫat či da, de bānde pohege?} \\
I mean, this writing in \ili{Pashto}, do you understand it? 

\speaker{S}\textit{Da ze wele nišama.} \\
`I cannot read that.' 
\myattrib{Int 1\_S, 191--194}
\end{interview}

The mother provides a more comprehensive reflection on the language use within her family and with respect to S. The children speak \ili{Pashto} with the parents and among themselves. In contrast to her older sisters, her mother reports that S has some difficulties in language learning. After the introduction of \ili{German}, these difficulties got even more important:

\begin{interview}\label{ch2.5.ex.18}
\speaker{MS}\textit{Aw os če peke ālmānee raḡla źma pe ḵyal če wāṛ'a ṭol gaḍ o waḍ ṣu. Koma ḵabara kawyi, neema pake ālmāni ṣeereka kee, źe ṣ Puҳto ṣeereka kee ter ḵpela ḵabara m ṣ āta oway...} \\
`Now \ili{German} is added to the mix as well, everything is getting mixed up completely…She says something, half of it is in \ili{German} and the other half is in \ili{Pashto}, that's how she communicates her sentences to me.' 
\myattrib{Int 1\_S\_mother, 170--173}
\end{interview}

Additionally, the mother assesses S's proficiency\is{Language!proficiency} in \ili{Pashto} as not very strong. That is why the mother tolerates her mixed use of \ili{German} and \ili{Pashto} in the family.

The mother speaks \ili{Urdu}, \ili{English}, and some \ili{Persian} (\ili{Dari}), and she mentions that her proficiency\is{Language!proficiency} in \ili{Persian} is improving, particularly in comprehension:

\begin{interview}\label{ch2.5.ex.19}
\speaker{I}\textit{Nor pa kumo žebo puyegey?} \\
`What other languages do you speak?' 

\speaker{MS}\textit{Puҳto, aw źam Parsi leka hos źe harźog oway kedayṣee jawab ṣ darkawelayṣom, kedayṣ ṣ ee pe ṣ ḍer na, ḵo puyegom pa ṭola.} \\
`I speak \ili{Pashto}, I also speak \ili{Persian} … but sometimes it takes a while for me to respond if you were to talk to me, I might not be able to respond if it's too much, but I understand \ili{Persian} perfectly.' 
\myattrib{Int 1\_S\_mother, 131--135}
\end{interview}

In the second interview with the child, which is still conducted in \ili{Pashto}, she declares \ili{German} her favorite language:

\begin{interview}\label{ch2.5.ex.20}
\speaker{I}\textit{Pašto? Ok. Te pa koma žaba ḫušala ye, ya pa koma žaba ḫabare kawal sta ḫwaḫigi?} \\
`\ili{Pashto}? Ok. In which language do you feel comfortable, which language do you like to speak most?' 

\speaker{S}\textit{Zama Deutsch bande ḫabare kawal ḫwaḫigi.} \\
`I like most to speak \ili{German} .' 
\myattrib{Int 2\_S 21--24}
\end{interview}

At the same time, she reports that \ili{Pashto} is still often used in the family and with other children at school who do not speak much \ili{German} (Int 2\_S, 46--52).

\subsubsection{E: Pashto and Dari}\label{ch2.5.ss.4.2.5}

The last case represents another constellation of languages, in Afghanistan and Germany, as well as some kind of openness with respect to different ethnic groups. The boy speaks \ili{Pashto} (and some \ili{German}) at home, but he also learned some \ili{Persian} with the neighbors' children in Germany:

\begin{interview}\label{ch2.5.ex.21}
\speaker{I}\textit{Noro žebo pege? Kum žebe bānde pege ḵbela de?} \\
`Do you speak any other languages? Which languages do you speak?' 

\speaker{E}\textit{Paҳto bānde... aw pārsi bānde, aw „bisschen“ deucṛ'} \\
`\ili{Pashto}... and \ili{Persian}... and some \ili{German}' 

\speaker{I}\textit{Hamaḡese pohege, pārsi de le kum jāi na zda keṛ'ai?} \\
`Where did you learn how to speak \ili{Persian}?' 

\speaker{E}\textit{Pārsi da dwa hālekāna di lande, hoḡo na me zda keṛ'i} \\
`I learned \ili{Persian} from the two boys down here' 
\myattrib{Int 1\_E, 29--46}
\end{interview}

The multilingual situation of Afghanistan, which seems to show, at least in this case, no hierarchy of the two official languages, finds a continuation even under the new conditions of life in Germany. The father gives additional information in this respect: 

In the beginning of the interview, E's father affirms: \textit{Mung moraney žeba paҳto da} `Our mother tongue is Pashto' (E\_father, 178f.). In the following, he refers also to the second common language in Afghanistan, \ili{Dari}:

\begin{interview}\label{ch2.5.ex.22}
\speaker{FE}\textit{Ho, māṣumān pa haḡa žeba ham ḵabere kai, aw bel māṣ ṣ umān, pe ṣ pārsi ham pohegi, pa pārsi dase pohegi če mesāl menga hamsāyagān be afḡānestān ke, awḡānestān ke ḵo ṭol ḵalk „gemiṣt“ žwand kai, Pārsiwānān, ṣ Paҳtāne „gemischt“ žwand kai.} \\
`Yes, the children also speak this language, and they also understand \ili{Persian}. \ili{Persian}, so they understand \ili{Persian} because, for example, our neighbours in Afghanistan. In Afghanistan all the people are mixed, the Persians, the Pashtuns, all mixed.' 
\myattrib{Int 1\_E\_father, 184--191}
\end{interview}

E's father differentiates between the two languages: \ili{Pashto} is their  heritage language, and the children speak \ili{Pashto} also with mother and father and among themselves (E\_father, 247--253). But \ili{Persian} (\ili{Dari}) is understood as well, due to the cultural and ethnic blending of Afghanistan. In their homeland, too, there was a mix of Persians (\textit{Pārsiwānān}) and Pashtuns (\textit{Paҳtāne}). This statement points to the ethnic diversity within Afghanistan, which proved to be relevant in the context of language and cultural integration. His son can profit from this in his hostel in Germany, where he can understand the \ili{Dari}-speaking boys.

In the second interview, which is still conducted in \ili{Pashto} with the boy, he claims that he feels comfortable with both languages, \ili{Pashto} and \ili{German}. As the family moved to another district, he still has no contact with other children outside of the school where he only speaks \ili{German} (Int 2\_E 24--34, 45--48).

The father describes more or less the same situation at home as in the first interview: The children speak more \ili{Pashto} among each other, while he tries to strengthen the use of \ili{German} also at home. Even \ili{Persian} is still used with other people from Afghanistan (Int 2\_E\_father, 51--83).

\subsection{Cultural identity and attitudes towards family languages}\label{ch2.5.ss.4.3}

\subsubsection{Z and her mother}\label{ch2.5.ss.4.3.1}

The second category is about language attitudes\is{Language!attitudes} and cultural identity. In the end of the interview with Z, she gives the following statement on \ili{Kurdish}:

\begin{interview}\label{ch2.5.ex.23}
\speaker{I}\textit{Was würdest du noch sonst, äh, gerne erzählen? Gibt's etwas, was du mir erzählen möchtest?} \\
`What else would you like to tell me? Is there something you would like to tell me?' 

\speaker{Z}\textit{Also zum Beispiel, äh, in viele Länder, also in Syrien, gibt's keine kurdische Sprache, die müssen in der Schule Arabisch lernen…. Und zum Beispiel meine Geschwister können deswegen nicht so viel Kurdisch. … Und ich denk', vielleicht wäre gut, wenn auch kurdische Schule gegeben wäre, dann konnte ich mehr über Kurdisch wissen … und kurdische Kultur und so.} \\
`For example, in many countries, in Syria, there is no \ili{Kurdish} language, they have to learn \ili{Arabic} at school… and, for example, my brothers and sisters don't know much \ili{Kurdish} because of that. … And I think maybe it would be good if there was also a \ili{Kurdish} school, then I could know more about \ili{Kurdish} … and \ili{Kurdish} culture and so about it.' 

\speaker{I}\textit{Ach so, über die Sprache}\\ 
`Oh, about the language' 

\speaker{Z}\textit{Mhm.} (approving)

\speaker{I}\textit{Dass diese Schule auf kurdischer Sprache ist?} \\
`That this school is in the \ili{Kurdish} language?' 

\speaker{Z}\textit{Ja, also nur für Kurden.} \\
`Yes, so only for the Kurds.' 
\myattrib{Int 1\_Z, 320--349}
\end{interview}

Z argues in favour of the teaching of \ili{Kurdish} at schools in Germany. As a reason for this, she refers to the exclusion of \ili{Kurdish} as school language in Syria. The missing \ili{Kurdish} teaching in the Syrian school is clearly seen in connection with tendencies of language loss in the case of her smaller siblings.

Z's mother, at the end of her interview, as well, expresses her gratitude with respect to our interest in the heritage language \ili{Kurdish}:

\begin{interview}\label{ch2.5.ex.24}
\speaker{I}And my last question, do you have any question for me? (asked in German)

(Z translates the question and the mother answers in Kurmanji\il{Kurdish!Kurmanji})

\speaker{MZ}\textit{Na, ez saet xweşbîn dikim. Bibê saet xweş. Bibê am hatine bîrê, em Kurmanc in. Em zimanê xwe nizanın û me ra deng kir.} \\
`No, I want to thank you. Thank you for thinking of us Kurds, too. We have no way of speaking our language, but she (MK) spoke to us about it.' 
\myattrib{Int 1\_Z\_mother, 569--576}
\end{interview}

Her mother repeatedly emphasises that it is important to her that her children can speak, read and write \ili{Kurdish}, as it is their language. But she points out that she cannot read and write and has difficulties teaching her children to do so in \ili{Kurdish}. Nontheless, she takes a proactive stance, which is expressed as a desire for a ``\ili{Kurdish} school'' in Germany: The school is attributed an important role in conveying knowledge about \ili{Kurdish} language and culture.

In the second interview, the girl identifies herself explicitly with the \ili{Kurdish} people:

\begin{interview}\label{ch2.5.ex.25}
\speaker{Z}\textit{JAA: Zum Beispiel wir Kurden. Es gibt sehr viele Kurden auf der Welt, in verschiedenen Ländern, zum Beispiel in der Türkei, Syrien, Irak und die haben halt kein eigene Land und zum Beispiel isch wurde auch in der Schule, wenn isch sage, isch komme also isch bin Kurde, sie sagen immer Fantasialand zu mir. Ja.} \\
`Yes, for example we Kurds. There are many Kurds in the world, in different countries, for example in Turkey, Syria, Iraq, and they do not have a country of their own. For example, in school, when I say, I come from… I'm \ili{Kurdish}, they say ``fantasyland'' to me. Yes. 

\speaker{I}\textit{Warum?} \\
`Why?' 

\speaker{Z}\textit{Ja, weil der Kurde keine eigene Land und keine eigene Sprache haben.} \\
`Yes, because the Kurds do not have a country and a language of their own.' 
\myattrib{Int 2\_Z 212--217}
\end{interview}

The girl reports an incident at school where she was derided when saying that she was \ili{Kurdish}. The expression \textit{Fantasialand} `fantasyland' in \ili{German} is used to refer to a place that exists only in fantasy, like a fairy tale. In this view, ethnic identity is linked to a state, without a state, no ethnic identity is possible. Z describes such a behavior as discriminating. 

When asked about her heritage language in the second interview, the mother admits the importance of \ili{German} for the future in Germany, even at the expense of the language she loves:

\begin{interview}\label{ch2.5.ex.26}
\speaker{I}\textit{Wie wichtig ist für Sie, dass Ihr Kind die Herkunftssprache in diesem Fall Kurdisch weiterspricht? Wie wichtig ist das für Sie?} \\
`How important is it for you that your child still speaks your language of origin, \ili{Kurdish} in this case? How important is that for you?' 

(Z translates the question. Z's mother answers in German.)

\speaker{MZ}\textit{Ja isch liebe meine Sprache ama hier alles Deutsch muss Deutsch lernen lesen, wir müssen lesen Deutsch} \\
`Yes, I love my language, but here everything is \ili{German}, we must learn to read \ili{German}, we must read \ili{German}' 
\myattrib{Int 2\_Z\_mother, 287--291}
\end{interview}

The final statement of the interview underlines again the importance of ``remembering our language'':

\begin{interview}\label{ch2.5.ex.27}
\speaker{MZ}\textit{Çiçikê em ser zimanê xwe deng dikin, em xwe dihesin!} \\
`We are talking about our language a bit, we remember ourselves! 

\speaker{Z}\textit{Sie findet es gut, weil sie redet so bisschen über ihre Sprache mit jemand anderem, also das findet sie gut.} (translates)\\
`She is happy because she can speak about her language with somebody else, she likes this.' 
\myattrib{Int 2\_Z\_mother, 592--595}
\end{interview}

\subsubsection{N's father}\label{ch2.5.ss.4.3.2}

N's father claims the same right for Kurds to maintain their own language, even reading and writing: 

\begin{interview}\label{ch2.5.ex.28}
\speaker{I}\textit{Zimanê we jî bo we girînge? Zimanê xwe jî hîn bibin, bixûnin, binivîsînin.} \\
`Is your language important for you? To learn, read and write your language?' 

\speaker{FN}\textit{Ez ci bibêjim. Zimanê me jî mihume, li vî derê em dixwazin li vî zimanê jî hîn bibin. Nivîsandinê, xwendinê. Em jî xwe re li vî derê karekî bikin.} \\
`How can I say this? Our language is important. We want them to learn our language here. To write, to read. And we can work here.' 
\myattrib{Int 1\_N\_parents, 436--440}
\end{interview}

After being asked about the importance of the language for his children, he concludes with the following words: 

\begin{interview}\label{ch2.5.ex.29}
\speaker{FN}\textit{Bo cinsîyeta wan, ji wan re jî muhîme. Ma her kes zimanê xwe bîrnake. Bila zimanê xwe, yê dayîkê bîr nekin.} \\
`It is important for their \isi{nationality}. Nobody forgets his language. One should not forget his language, his mother tongue' 
\myattrib{Int 1\_N\_parents, 517--523}
\end{interview}

The father emphasizes the importance of the heritage language even for his children in the new country. As a reason for this, he refers to the children's \isi{nationality} (\textit{cinsîyeta wan}). In the last two sentences, he moves from stating a mere fact to fixing a moral obligation. 

In the second interview, the importance of the heritage language is again underlined, but the teaching of \ili{Kurdish} is not in focus:

\begin{interview}\label{ch2.5.ex.30}
\speaker{I}\textit{Dibêje, äh N, keça te, dibêje, tere di medresê da zimanê Kurdî, yanî yanî mh heye dersê wê tişt?} (translates)\\
`She (2nd interviewer) asks, if N, your daughter, goes to \ili{Kurdish} lessons at school. Do such lessons exist?'

\speaker{FN}\textit{Na, na, bes bi Almanî.} \\
`No, only in \ili{German}' 

\speaker{I}\textit{Bi Almanî, Kurdî nîne?} \\
`In \ili{German}, not in \ili{Kurdish}? 

\speaker{FN}\textit{Na} \\
`No' 
\myattrib{Int 2\_N\_parents, 303--310}
\end{interview}

The girl agrees that she would like to read and write in \ili{Kurdish}, but later she speaks about \ili{Arabic}, which is practiced in a mosque:

\begin{interview}\label{ch2.5.ex.31}
\speaker{I}\textit{Ja. Okay. Und möchtest du gerne deine kurdische Sprache lesen und schreiben lernen?} \\
`Do you want to learn to read and write your \ili{Kurdish} language?' 

\speaker{N}\textit{Ja} \\
`Yes' 

\speaker{I}\textit{Möchtest du das? Okay.} \\
Do you want that? Okay.

\speaker{N}\textit{Also, gibt gar nicht, wenn wir lesen auf Kurdisch, sondern gibt auf Arabisch.} \\
`Reading in \ili{Kurdish} does not exist, it's in \ili{Arabic}' 

\speaker{I}\textit{Aha diese Kurdisch ist auch arabische Schrift. Meinst du, oder? Kannst du Arabisch lesen?} \\
`This \ili{Kurdish} is in \ili{Arabic} \isi{script}, isn't it? Can you read \ili{Arabic} \isi{script}?' 

\speaker{N}\textit{Nein, aber ich lerne.} \\
`No, but I learn.' 
\myattrib{Int 2\_N 10:21--10:36}
\end{interview}

N denies here that \ili{Kurdish} can be written, but claims that only \ili{Arabic} has a \isi{script}. She is supposed to learn this \isi{script} in the mosque, but up till now they only participate in the prayer:

\begin{interview}\label{ch2.5.ex.32}
\speaker{N}\textit{So eine Moschee und wir beten da. Aber vielleicht können wir da auch lernen.} \\
`Such a mosque and we pray there. But perhaps we can also learn there.' 
\myattrib{Int 2\_N 11:14}
\end{interview}
 
\subsubsection{A's mother}\label{ch2.5.ss.4.3.3}

The mother of the third \ili{Kurdish} student also wants to maintain the \ili{Kurdish} language in the family. 

\begin{interview}\label{ch2.5.ex.33}
\speaker{I}\textit{Wie wichtig ist für Sie, dass Ihre Kinder Bahdini\il{Kurdish!Bahdini} lernen oder Ihre Sprache lernen? Und lesen, schreiben?} \\
`How important is it for you that your children learn Bahdini\il{Kurdish!Bahdini} or your language? And read, write?' 

\speaker{A}\textit{Çima girîng e ku em fêrî axaftin û nivîsandina Kurdî bibin?}  (translates)\\
`Why is it important that that we learn to speak and write \ili{Kurdish}?' 

\speaker{MA}\textit{Divê em Kurdiya xwe ji bîr nekin.} \\
`We should not forget our \ili{Kurdish}.' 

\speaker{A}\textit{Sie will nicht, dass wir Kurdisch vergessen, Ja.} (translates)\\
`She doesn't want that we forget \ili{Kurdish}.' 

\speaker{I}\textit{Das ist Ihnen sehr wichtig, Ihre Sprache?} \\
`This is very important for you, your language?' 

\speaker{MA}\textit{Ja} \\
`Yes' 
\myattrib{Int 1\_A\_mother, 359--365}
\end{interview}

The mother of the third \ili{Kurdish} student also wants to maintain the \ili{Kurdish} language in the family. 

\begin{interview}%\label{ch2.5.ex.33}
\speaker{I}\textit{Wie wichtig ist für Sie, dass Ihre Kinder Bahdini\il{Kurdish!Bahdini} lernen oder Ihre Sprache lernen? Und lesen, schreiben?} \\
`How important is it for you that your children learn Bahdini\il{Kurdish!Bahdini} or your language? And read, write?' 

\speaker{A}\textit{Çima girîng e ku em fêrî axaftin û nivîsandina Kurdî bibin?} (translates)\\
`Why is it important that that we learn to speak and write \ili{Kurdish}?' 

\speaker{MA}\textit{Divê em Kurdiya xwe ji bîr nekin.} \\
`We should not forget our \ili{Kurdish}.' 

\speaker{A}\textit{Sie will nicht, dass wir Kurdisch vergessen, Ja.} (translates)\\
`She doesn't want that we forget \ili{Kurdish}.' 

\speaker{I}\textit{Das ist Ihnen sehr wichtig, Ihre Sprache?} \\
`This is very important for you, your language?' 

\speaker{MA}\textit{Ja.} \\
`Yes.' 
\myattrib{Int 1\_A\_mother, 359--365}
\end{interview}

The mother emphasizes the importance of not forgetting \ili{Kurdish}. While the official language in Iraq is \ili{Arabic} and is needed in the administration, \ili{Kurdish} as a spoken language is more relevant in the family and sociocultural sphere. Although the interview question was translated very precisely into \ili{Kurdish} here (\textit{ku em fêrî axaftin û nivîsandina kurdî bibin} `that we learn to read and write \ili{Kurdish}'), the mother only gives a very general answer: the language should not be forgotten. This family takes a more defensive stance towards the question of literacy in the family language \ili{Kurdish}. She does not express specific wishes or measures how the language could be preserved in Germany. 

\subsubsection{S's mother (Pashto)}\label{ch2.5.ss.4.3.4}

In the case of \ili{Pashto}, the main argument of the mother is the necessity to be able to communicate with their family. While the question focuses the official language, the answer relates to the family language. She underscores the cultural importance of the language, as it not only serves as a means of communication, but also strengthens the cultural identity and connection of the family. It is evident that the mother appreciates the conveying of the language to the next generation in order to preserve cultural and familial bonds:

\begin{interview}\label{ch2.5.ex.34}
\speaker{I}\textit{Stāso de asle hewād rasmi žeba źe ahmyiat lari?} \\
`What significance does your official language have in your home country, how important is it for you?' \\
\speaker{MS}\textit{ḍer mohema da... ahmiyat lari} \\
`Very important... very important part [it plays]' 

\speaker{I}\textit{Stāso māṣum da para rasmi žeba ze ma´nā lari? Yāni māṣ ṣ um tam źe ahmyiat lari ṣ} (translates)\\
`What role does the official language [of their home country] play for your children, how important is this language for them?' 

\speaker{MS}\textit{ham dumra ahmiyat lari no... oḡoy da para zarur di, zeka če ḵpel famil sara,  ḵpel koraney sara munga če yaw zay patikigo bayād če hoḡoyla puҳto warsi.} \\
`It's exactly the same important thing... They need it because when they are with their family [\textit{fāmil}], together with their family [\textit{koraney}], they have to be able to speak \ili{Pashto}.' 
\myattrib{Int 1\_S\_mother, 110--121}
\end{interview}

Learning to read and write \ili{Pashto} in Germany, however, is not an issue here. In the second interview, the mother makes this point even more explicit:

\begin{interview}\label{ch2.5.ex.35}
\speaker{MS}\textit{Moranәy žaba hagha weyli ši, ḫo dase lustale ba na ši laka či za wem, dase hagha lustale na ši. Zaka če waṛa da.} \\
`She can express herself in the mother tongue, but she will not be able to read it as I do. She is not prepared for that, because she is still young.' 
\myattrib{Int 2\_S\_mother, 171--173}
\end{interview}

\subsubsection{E's father (Pashto)}\label{ch2.5.ss.4.3.5}

When talking about language use at home, the father claims to stimulate the use of \ili{German}, as well:

\begin{interview}\label{ch2.5.ex.36}
\speaker{FE}\textit{Paҳto bānde ḵabere kawi... ze pe de waḵto ke worta lagiya yem, māle tāso pa ḵpel mābayn ke pa deuč ḵabere kawey če tāso deuč ҳa ṣṣi, paҳto preda, awḡānestān pāte ṣṣu, dalta ḵo če deuč de neyi zda, ḍer muṣṣkelāt di.} \\
`In \ili{Pashto}... I try to say to them always, that they need to speak \ili{German} among each other, leave \ili{Pashto} at the side, Afghanistan is far away, if you don't speak \ili{German} here, you will have many problems.' 
\myattrib{Int 1\_E\_father, 253--259}
\end{interview}

Already at this stage of the interview, the importance of \ili{German} is underlined in order to avoid problems, even if the children use it among each other. In the same vain, the father admonishes the children not to watch Afghan movies, but \ili{German} television programs:

\begin{interview}\label{ch2.5.ex.37}
\speaker{FE}\textit{ze zalerwist sāta dā ta´qiq pe kawma, če ze pe kār ke yem, duy ta ṭilifun kawm, ze worta waym če bel flm, afḡānay flm, yā bel ṣṣay bel ṣṣay, dā ṣṣay merawarey, be ḡair le deuč na.} \\
In the first place at school, in the second at home. I'm after them 24 hours a day, even from work I ring them up and say to them, that they should not watch other or Afghan movies, but only \ili{German}.' 
\myattrib{Int 1\_E\_father, 272--279}
\end{interview}

The father comes back later to the importance of Afghanistan's two official languages, which are no longer important in Germany: 

\begin{interview}\label{ch2.5.ex.38}
\speaker{FE}\textit{Moraney žeba moraney žeba i role daḡa wa če pe ḵpel watan ki žwand kai, bāyad či dā dwaṛ'e de zda yi, kana? Aw delta raḡlu, delta mung takum ahmiyat ne lari, delta mung ta mohema deuč da.}\\ `The mother tongue is the mother tongue, so the role that it plays it plays when you live in your home country, you have to speak both languages there, right? But here, here it doesn't play a role for us, here is \ili{German} important for us.'
\myattrib{Int 1\_E\_father, 438--440} 
\end{interview}

With such a tautological statement, the father reproduces the ideology\is{Language!Ideology} of the mother tongue, but calls it into question at the same time. The situation in the country of origin where \ili{Pashto} and \ili{Dari} are important, is juxtaposed to Germany, where only \ili{German} is still important.

\begin{interview}\label{ch2.5.ex.39}
\speaker{I}\textit{stāso da pāra źumra mohema da ter źo stāso māṣṣumān ālmāni zda kṛ’i?}
`How important is it for you that your child learns \ili{German}?' 

\speaker{FE}\textit{Mung ta, le har źe na ḍera mohema da, biḵi de de de mohemyat źe nesem derta welay, zeҳta ḍera rāta mohema da če de mā māṣumān ālmāni ṣ žeba zda kṛ'i.} \\
`For us it is more important than anything else, I cannot exactly define the importance of this right now, but it is extremely important for me that my children learn learn \ili{German}. 
\myattrib{Int 1\_E\_father, 450--455}
\end{interview}

In such a way, the unconditional importance of \ili{German} and the learning of \ili{German} is again underlined in the final passage of the interview. The father makes a direct connection between language and culture when asked about the importance of the \ili{German} language:

\begin{interview}\label{ch2.5.ex.40}
\speaker{FE}\textit{Pa ālmān ke žwand kaw, demung ḵo no nor ham delta žwand da no, bāyad de de ḵalko kultur, kānun, har či ta ehterām woku.} \\
We live in Germany now, our life is here, so we must also respect the culture and laws here.' 
\myattrib{Int 1\_E\_father, 388--391}
\end{interview}

This statement could be understood in terms of an obligation to integration; in any case, instrumental learning motivation is not the primary focus here. In the second interview, this strong preference of \ili{German} is to some point counterbalanced by a more favorable view on the heritage language:

\begin{interview}\label{ch2.5.ex.41}
\speaker{I}\textit{Stāso lapāra dā tsomra mohima de če stāso māšom pa moranәy žaba ḫabare kawal ta dawām warkṛi?}  \\
`How important is it for you that your children learn the heritage language?'

\speaker{FE}\textit{Monga ghwaṛu če dwaṛa izda wi, paḫto ao ham deutsch.} \\
`We want them to speak both languages, \ili{Pashto} and \ili{German}.' 

\speaker{I}\textit{Ao dā stāso lapāra dā tsomra mohima de če stāso māšom hum pa moranәy žaba lostal ao likal kawale ši?}\\ 
`How important is it for you that your children can read and write the mother tongue?' 

\speaker{FE}\textit{Da dwaṛa mohima di bayt če moranәy žaba di.} \\
`Both are very important. It is obligatory to know the mother tongue.' 
\myattrib{Int 2\_E\_father 6, 152--161}
\end{interview}

\subsection{Typology of the case studies}\label{ch2.5.ss.4.4}

The language policies\is{Language!policies} of the five families can be described in the theoretical framework of \citet{darvin_norton2015,darvin_norton2017} as follows. 

\subsubsection{Linguistic and cultural capital}

The children's first \isi{minority} language learnt from early childhood onwards, which was \ili{Kurdish} in three and \ili{Pashto} in two cases, is complemented quite early by other languages like \ili{Arabic} (official language in Syria and Iraq) and \ili{Dari} (second official language in Afghanistan). Only two older girls (Z and A) have school experience in their country of origin, but due to the unstable political situation, the knowledge of the written language got lost. While the speakers of \ili{Kurdish} emphasize the difference between their family language and the official national language, \ili{Arabic}, which is no longer important for them and their children, the \ili{Pashto}-speaking Afghan families have a more tolerant view of \ili{Dari}, Afghanistan's second official language. 

The younger children also had some access to other languages of the country of origin (\ili{Dari} in case E), but often the stay in a third country (like Turkey, Greece or Pakistan) had a greater impact. In Germany, \ili{Turkish} still plays an important role for Z, while E still uses \ili{Dari} with some friends. The status of the heritage language, still quite stable in early childhood, begins to lose ground, especially in the case of younger siblings (like S, herself a younger sibling, or the siblings of Z and A).

A typical characteristic of many migrant families is the focus on the language of origin in conversation with parents and a greater use of the language of the immigration country among children and their peers \citep{portes_fernandez-kelly_haller2009,kempert_et_al2016}. In our data, however, some of the children claim to speak L1 almost exclusively inside of the family. In other interviews, the use of \ili{German} at home is even encouraged by the parents, such as E's father, who integrates \ili{German} into his conversations with the children in order to help them to learn it.

\subsubsection{Language ideologies}

The beliefs of parents and children on their linguistic situation are quite heterogeneous, at least in the first interviews. Relying on different typologies of sociolinguistic orientations of \isi{minorities} in migration contexts \citep{berry1997,brizic2021,brizic_bulut_simsek2021}, we differentiated three stances, that the families take with respect to the future of their heritage languages: 

\begin{itemize}
        \item Most parents and children consider their heritage language to be a part of their cultural identity that they should maintain (cases 2, 3, 4: Iraq, Syria, Afghanistan). This stance in favor of the heritage language remains defensive: an ideology of cultural identity is claimed, but no investments are proposed. For mother S, learning \ili{German} represents the dominant reality and \ili{Pashto} is only a wish. A's mother and N's father repeat the formula of not forgetting the language without taking measures for literacy learning. The stance corresponds to Brizić's ``intermediate voice'' \citep{brizic2021}.
        \item Only Z's family wishes to have a school teaching in the \ili{Kurdish} \isi{minority} language (case 1: Kurdistan/Syria). This proactive stance in favor of the heritage language underscores the necessity of investment in order to maintain the language, combined with an idea of dual integration. Learning to read and write in this language is not anymore a mere formula. In Brizić's terms, this is the ``balanced voice''.
%
        \item Some parents fear that maintaining or fostering the heritage language could impede the successful acquisition\is{Language!acquisition} of the \ili{German} language or would not be beneficial to the child. In this case, the father opts for a strong investment in the learning of \ili{German}, combined with an ideology of assimilation, the ``targeting voice'' in Brizić's study (case 5: Afghanistan).
%
\end{itemize}
    
The second interviews show some homogenizing processes: In all cases, the idea of maintaining the heritage language gets more important (even case 5), while the proactive stance of case 1 loses ground under the pressure of \ili{German} and the feeling of discrimination against \ili{Kurdish} identity.
    
\subsubsection{Investment}
While all families adapt themselves to a future in Germany and support their children in the learning process of \ili{German}, the importance of communicating with the greater family in the country of origin still plays a role. Thus, in most cases, oral communication in the family language still has a social and emotional function, while they do not attribute to their dominant heritage language (\ili{Kurdish} and \ili{Pashto}) an ``official'' status as a written language. The oral use of the language in conversation with the family and people from the country of origin is considered sufficient for the maintenance\is{Language!maintenance} of a \ili{Kurdish} or Afghan identity.
    
Community support, discrimination, and educational opportunities were treated only indirectly in the interviews. In all cases, parents deny membership in or a closer relationship to associations of their migrant community. Contacts with compatriots exist in the hostels in which newly arrived families live. Parents avoid mentioning discrimination, even if they indicate few contacts with people of \ili{German} origin in their neighborhood. In the cases of newly arrived children (N, E and S), the opportunities for contact with \ili{German}-speaking children at school, which were described as low in the first interviews, get more important in the second interview when these children were integrated in regular classes. 

\section{Summary}\label{ch2.5.ss.5}

The analysis of language attitudes\is{Language!attitudes} and family language policies\is{Language!policies} in newly arrived families of \ili{Kurdish} and \ili{Pashto} background in Germany during the first two years of their stay shows the importance of several languages: in addition to a \isi{minority} language spoken in the family we find influence of the official languages of the country of origin and of transit languages of countries where some families stayed for a longer time. Conflicting ideologies of learning \ili{German} as fast as possible and maintaining the heritage language in the family play a role in all families. But while investment in the \ili{German} school is considered essential by all families, decisive steps to prevent the loss of the heritage language are rare.

Some limitations of this study should be mentioned: In order to establish contact with newly arrived immigrants as soon as possible, the \ili{German} schools served as the starting point of the project. Therefore, only a small sample of cases with very divergent linguistic and cultural backgrounds could be included. A commonality, however, is a quite late starting point in second language acquisition\is{Language!acquisition}, similar age, forced migration, and the \isi{minority} status in the country of origin. The cases are by no means representative of their whole community. For speakers of \ili{Kurdish}, similar options exist in the literature on other communities \citep{brizic2021}. For speakers of \ili{Pashto}, to date, research on language acquisition\is{Language!acquisition} and attitudes\is{Language!attitudes} does not exist (but see for \ili{Dari}-speaking families from Afghanistan: \citealt{panagiotopoulou_uca2023_dynamic,panagiotopoulou_uca_samani2023}).

In summary, this study contributes to our understanding of the complex dynamics surrounding \isi{minority} language retention within migrant families in Germany. The diversity in the interplay between language, cultural identity, and integration, as reflected in our case studies, might encourage educators to take into account the different aspirations of children and their parents. The introduction of \isi{minority} languages as a subject in the \ili{German} school, opted for in Z's case, may not be an adequate solution for E's family. In any case, more inclusive schooling practices may be helpful for all the children concerned.

\sloppy
\printbibliography[heading=subbibliography,notkeyword=this]
\end{document}
